Les 200 adhérents d'une salle de sport ne pratiquent qu'une seule activité parmi les deux activités suivantes : le step et le crossfit. La répartition des adhérents est donnée dans le tableau suivant.
	
	
	\begin{center}
		\begin{tabularx}{0.85\linewidth}{|*4{>{\vspace*{\stretch{1}}\centering\arraybackslash}X|}}
			\cline{2-4}
			\multicolumn{1}{c|}{}		& Step & Crossfit & Total \\[0.5cm]
			\hline		
			Homme   & 20   & 80       & 100   \\[0.5cm]
			\hline
			Femme   & 60   & 40       & 100   \\[0.5cm]
			\hline
			Total   & 80   & 120      & 200   \\[0.5cm]
			\hline
		\end{tabularx}
	\end{center}
	
	On choisit un adhérent au hasard parmi les 200 adhérents.
	
	On considère les évènements suivants :
	
	\begin{itemize}[left=5mm,label={\textbullet~~}]
		\item $F$ : \og l'adhérent est une femme \fg{} ;
		\item $H$ : \og l'adhérent est un homme \fg{} ;
		\item $S$ : \og l'adhérent pratique le step \fg{} ;
		\item $C$ : \og l'adhérent pratique le crossfit \fg{}.
	\end{itemize}
	
	\begin{enumerate}
		\item Déterminer la probabilité $\text{P}(F)$ de l'évènement $F$.
		
		\item Déterminer la probabilité que l'adhérent soit un homme qui pratique le step.
		
		\item Déterminer la probabilité de l'évènement $F\cap S$.
		
		\item Les évènements $F$ et $S$ sont-ils indépendants ? Justifier la réponse.
		
		\item On choisit au hasard une femme parmi les adhérents.
		
		Quelle est la probabilité qu'elle pratique le crossfit ?
		
		\item Déterminer la probabilité $\text{P}_{C}(F)$.
	\end{enumerate}
